\chapter{Zadatak 4a --- jesu li modovi Hurwitzovi}
\label{ch:z4-hurwitz}

\section{Što se ovdje traži}

Izračunati svojstvene vrijednosti obiju matrica i reći je li svaki mod zasebno
stabilan.

\begin{zadatak}{Zadatak 4, podzadatak a)}
a) Jesu li matrice $A_{\mathrm{parno}}$ i $A_{\mathrm{neparno}}$ stabilne matrice
(Hurwitzove matrice)? Odnosno, jesu li oba dinamička moda stabilna?
\emph{(Napomena: ovo možete provjeriti izračunom svojstvenih vrijednosti ovih
matrica.)}
\end{zadatak}

\begin{predznanje}
\begin{itemize}[nosep,leftmargin=*]
  \item Karakteristična jednadžba $2\times2$ matrice --- \S\ref{sec:svojstvene}.
  \item Što znači Hurwitz --- \S\ref{sec:oznake-dinamika}.
  \item Trag i determinanta --- \S\ref{sec:vektori-matrice}.
\end{itemize}
\end{predznanje}

\begin{pojam}{Hurwitzova matrica}
Kvadratna matrica $A$ je \textbf{Hurwitzova} ako sve njezine svojstvene
vrijednosti imaju strogo negativan realni dio, $\Real \lambda_i < 0$. Za linearan
vremenski invarijantan sustav $\dot{x} = Ax$ to je nužan i dovoljan uvjet
asimptotske (i eksponencijalne) stabilnosti: rješenje je
$x(t) = e^{At}x(0)$, a $e^{At} \to 0$ točno kad su svi $\Real \lambda_i < 0$.
\end{pojam}

\section{Izvod za parni mod}
\label{sec:z4-parni}

\begin{korak}[Karakteristična jednadžba, s predznacima]
\begin{equation}
A_{\mathrm{parno}} = \begin{bmatrix} -1 & -\tfrac{\pi}{6} \\[2pt] \tfrac{3\pi}{2} & -1 \end{bmatrix}
\;\Longrightarrow\;
\lambda I - A_{\mathrm{parno}} = \begin{bmatrix} \lambda + 1 & \tfrac{\pi}{6} \\[2pt] -\tfrac{3\pi}{2} & \lambda + 1 \end{bmatrix}.
\end{equation}
Determinanta je umnožak dijagonale minus umnožak antidijagonale:
\begin{equation}
\det(\lambda I - A_{\mathrm{parno}})
= (\lambda+1)^2 - \left(\frac{\pi}{6}\right)\left(-\frac{3\pi}{2}\right)
= (\lambda+1)^2 + \frac{\pi^2}{4} = 0 .
\label{eq:z4-karakteristicna}
\end{equation}
Odatle
\begin{equation}
(\lambda + 1)^2 = -\frac{\pi^2}{4}
\;\Longrightarrow\;
\lambda + 1 = \pm j\frac{\pi}{2}
\;\Longrightarrow\;
\lambda_{1,2}(A_{\mathrm{parno}}) = -1 \pm j\frac{\pi}{2} = -1 \pm 1{,}570796\,j .
\end{equation}
\end{korak}

\begin{zamka}[Zamka: predznak umnoška izvandijagonalnih elemenata]
Ovdje je vrlo lako izgubiti predznak i zaključak i dalje ispadne točan --- pa se
greška ne primijeti. Pazi:
\begin{itemize}[nosep,leftmargin=*]
  \item Izvandijagonalni elementi \emph{same matrice} $A_{\mathrm{parno}}$ su
        $-\tfrac{\pi}{6}$ i $\tfrac{3\pi}{2}$; njihov umnožak je
        $\boldsymbol{-\tfrac{\pi^2}{4}}$, dakle \emph{negativan}.
  \item U karakterističnoj jednadžbi taj se umnožak \emph{oduzima}, pa
        $-(-\tfrac{\pi^2}{4}) = +\tfrac{\pi^2}{4}$ --- odatle plus u
        \eqref{eq:z4-karakteristicna}.
\end{itemize}
Reći ,,umnožak izvandijagonalnih elemenata je $+\pi^2/4$`` je netočno; točno je
da je \emph{determinanta} matrice $A$ jednaka
$\det A = (-1)(-1) - (-\tfrac{\pi}{6})(\tfrac{3\pi}{2}) = 1 + \tfrac{\pi^2}{4}$.
Rezultat je isti, ali formulacija je ono što ispitivač sluša.
\end{zamka}

\begin{korak}[Kontrola preko traga i determinante]
Za $2\times2$ matricu vrijedi $\lambda^2 - (\tr A)\lambda + \det A = 0$. Ovdje:
\begin{equation}
\tr A_{\mathrm{parno}} = -1 + (-1) = -2,
\qquad
\det A_{\mathrm{parno}} = 1 + \frac{\pi^2}{4} \approx 3{,}4674,
\end{equation}
pa je $\lambda^2 + 2\lambda + \bigl(1 + \tfrac{\pi^2}{4}\bigr) = 0$, što je
identično \eqref{eq:z4-karakteristicna}. Uz to:
\begin{itemize}[nosep,leftmargin=*]
  \item $\lambda_1 + \lambda_2 = \tr A = -2$, a doista $(-1+j\tfrac{\pi}{2}) + (-1-j\tfrac{\pi}{2}) = -2$; \checkmark
  \item $\lambda_1 \lambda_2 = \det A = 1 + \tfrac{\pi^2}{4}$, a doista
        $(-1)^2 + (\tfrac{\pi}{2})^2 = 1 + \tfrac{\pi^2}{4}$. \checkmark
\end{itemize}
\end{korak}

\section{Neparni mod: isti spektar}
\label{sec:z4-neparni}

\begin{korak}
\begin{equation}
A_{\mathrm{neparno}} = \begin{bmatrix} -1 & -\tfrac{3\pi}{2} \\[2pt] \tfrac{\pi}{6} & -1 \end{bmatrix}.
\end{equation}
Trag je opet $-2$, a determinanta
\begin{equation}
\det A_{\mathrm{neparno}} = 1 - \left(-\frac{3\pi}{2}\right)\left(\frac{\pi}{6}\right) = 1 + \frac{\pi^2}{4},
\end{equation}
dakle \emph{identična} onoj parnog moda --- jer su izvandijagonalni elementi samo
zamijenili mjesta, a umnožak je komutativan. Karakteristična jednadžba je zato
ista, pa je
\begin{equation}
\lambda_{1,2}(A_{\mathrm{neparno}}) = -1 \pm j\frac{\pi}{2}.
\end{equation}
\end{korak}

\begin{formula}{Odgovor na podzadatak a)}
\begin{equation}
\lambda_{1,2}(A_{\mathrm{parno}}) = \lambda_{1,2}(A_{\mathrm{neparno}}) = -1 \pm j\frac{\pi}{2} .
\end{equation}
Realni dio svih svojstvenih vrijednosti je $-1 < 0$, pa su \textbf{obje matrice
Hurwitzove} i svaki je mod \emph{zasebno} asimptotski stabilan.
\end{formula}

\begin{lstlisting}[style=outputstyle]
--- (a) Stabilnost modova ---
eig(A_parno)   = -1.000000 +1.570796i,  -1.000000 -1.570796i
eig(A_neparno) = -1.000000 +1.570796i,  -1.000000 -1.570796i
pi/2 = 1.570796 -> obje matrice su Hurwitzove (Re = -1 < 0)
\end{lstlisting}

\section{Što svojstvene vrijednosti fizikalno znače}
\label{sec:z4-fizikalno}

\begin{pojam}{Čitanje kompleksnog para $\sigma \pm j\omega$}
\begin{itemize}[leftmargin=*]
  \item \textbf{Realni dio $\sigma = -1$} je brzina prigušenja: amplituda pada
        kao $e^{-t}$, dakle na $37\%$ u jednoj sekundi.
  \item \textbf{Imaginarni dio $\omega = \pi/2 \approx 1{,}571$ rad/s} je kružna
        frekvencija titranja. Period je $T = 2\pi/\omega = 4$ s.
  \item \textbf{Kompleksan par} znači da odziv \emph{oscilira} dok se gasi. U
        prostoru stanja to je spirala prema ishodištu.
\end{itemize}
\end{pojam}

\begin{korak}[Zašto je $\omega = \pi/2$ ovdje presudan]
Period titranja svakog moda je $4$ s, a mod se drži $1$ s. U jednoj sekundi
trajektorija dakle prijeđe točno \emph{četvrtinu} kruga, odnosno zakrene se za
$90^\circ$. Za to vrijeme amplituda padne faktorom $e^{-1} \approx 0{,}368$.

Ta dva učinka --- rotacija za $90^\circ$ i prigušenje za $e^{-1}$ --- međusobno se
natječu. Podzadatak b) pokazat će da rotacija pobjeđuje: zakret od $90^\circ$
uzastopno pogađa smjer u kojem drugi mod ima veliko pojačanje, pa se stanje
između dva prespajanja pojača više nego što ga prigušenje uspije smanjiti.
\end{korak}

\begin{zamka}[Zamka: ,,oba su moda stabilna, dakle sustav je stabilan``]
Ovo je \emph{fatalna} zamka cijelog Zadatka 4, i cijeli podzadatak b) postoji da
je razotkrije. Hurwitzovost svakog moda zasebno \textbf{ne} povlači stabilnost
sklopnog sustava. Odgovor pod a) mora glasiti ,,oba su moda zasebno stabilna``,
i tu stati --- bez zaključka o sustavu.
\end{zamka}

\begin{izgovor}
,,Karakteristična jednadžba parnog moda: determinanta od lambda i minus a jednaka
je lambda plus jedan na kvadrat, minus umnožak izvandijagonalnih elemenata
matrice lambda i minus a. Ti elementi su pi kroz šest i minus tri pi kroz dva, pa
je njihov umnožak minus pi na kvadrat kroz četiri, i kad se oduzme, dobiva se
plus pi na kvadrat kroz četiri.

Dakle lambda plus jedan na kvadrat plus pi na kvadrat kroz četiri jednako nuli,
odakle lambda jednako minus jedan plus minus je pi kroz dva.

Za neparni mod izvandijagonalni elementi su samo zamijenili mjesta, umnožak je
isti, pa je i spektar isti.

Realni dio je minus jedan, dakle obje su matrice Hurwitzove i svaki je mod
zasebno asimptotski stabilan. Naglasio bih da iz toga \emph{ne} slijedi da je
sklopni sustav stabilan --- to provjeravam u sljedećem podzadatku.``
\end{izgovor}
